\documentclass{article}

% AI4Mat @ NeurIPS 2026 double-blind workshop submission.
% This manuscript is intended as a full-length submission.
\usepackage[dblblindworkshop]{neurips_2026}
\workshoptitle{AI for Accelerated Materials Design (AI4Mat)}

\usepackage[utf8]{inputenc}
\usepackage[T1]{fontenc}
\usepackage{amsmath,amssymb,mathtools}
\usepackage{booktabs}
\usepackage{graphicx}
\usepackage{xcolor}
\usepackage{microtype}
\usepackage{hyperref}
\usepackage{url}

\newcommand{\R}{\mathbb{R}}
\newcommand{\HV}{\operatorname{HV}}
\newcommand{\cf}{\mathrm{CF}}
\newcommand{\direct}{\mathrm{direct}}
\newcommand{\composemap}{\ell}
\newcommand{\resultplot}[2][]{%
  \includegraphics[#1]{#2}%
}

\title{Composite Multi-Objective Bayesian Optimization for Scientific Design: An Application Study}
\author{}

\begin{document}
\maketitle

\begin{abstract}
Scientific design seeks materials, formulations, and processes that balance
competing goals, such as improving performance while reducing cost or
environmental impact.  Finding these tradeoffs is difficult because experiments
and simulations are expensive, their input--output relationships are unknown,
and their design spaces can be high-dimensional.  Multi-objective Bayesian
optimization (MOBO) addresses this problem by using past evaluations to decide
what to test next.  However, a scientific evaluation often produces a rich
response---such as a spectrum, trajectory, or set of physical
contributions---before that response is reduced to a few final scores.  Most
MOBO methods model only those scores and discard the intermediate information.
Bayesian optimization of composite functions (BOCF) was originally developed
for scalar objectives.  We apply the same principle to multi-objective
scientific design: model the observed response, pass posterior samples through
the known scientific reduction, and use the resulting objective samples inside
an existing MOBO method.  This is an application study rather than a new
acquisition method.  We compare direct and composite representations across
four MOBO approaches on two synthetic and three scientific benchmarks, with
two objectives and input dimensions ranging from 4 to 600.  The suite is
defined by an equal-score/different-sensitivity criterion: two distinct
intermediate responses can receive the same final score while responding
differently to a small component change.  Under matched evaluation budgets,
composite modeling achieves higher mean final hypervolume in all ten
method--benchmark comparisons, with a median relative gain of 6.98\%.
Preserving information already produced by an experiment can therefore improve
scientific search without requiring additional physical evaluations.
\end{abstract}

\section{Introduction}
\label{sec:introduction}

Scientific design is the task of choosing experimental or engineering
parameters that produce systems with desired properties.  It is central to
materials discovery, chemical formulation, and process optimization, and it is
a core function of autonomous laboratories that connect computational design,
automated synthesis, and characterization
\cite{lookman2019active,stach2021autonomous}.  These applications rarely have a
single goal.  A useful design may need to balance activity and stability,
optical response and manufacturability, or performance and environmental cost.
The scientific output is therefore a set of useful tradeoffs rather than one
best design.

The challenge is that scientific evaluations are expensive and the relation
between design choices and measured properties is usually unknown.  A
formulation experiment, for example, may vary several ingredient
concentrations, require physical preparation, and return a noisy spectrum only
after substantial laboratory time.  The number of possible formulations grows
quickly with the number of ingredients, while only a small fraction can be
tested.  Similar cost and dimensionality constraints occur in simulation-based
mechanics, ecosystem control, and materials processing.

Bayesian optimization addresses this setting by building a probabilistic model
from previous evaluations and using it to choose the next experiment
\cite{shahriari2016taking,frazier2018tutorial}.  Its multi-objective extension
(MOBO) seeks a Pareto set of designs for which no objective can be improved
without worsening another.  MOBO has consequently become important in
materials design and self-driving laboratories
\cite{gopakumar2018multiobjective,low2024constrained,macleod2022selfdriving}.

High-dimensional design spaces make this search harder because both the
predictive model and the selection of new experiments become less reliable.
The methods studied here represent two scalable responses.  Spherical-linear
models use a compact geometric representation, while MORBO coordinates several
local trust regions instead of fitting one global search model
\cite{doumont2026linear,daulton2022morbo}.

Scalarization provides a complementary way to handle several objectives.  It
turns different preferences among the objectives into a sequence of
single-objective searches.  Smooth Tchebycheff scalarization can therefore be
paired with both conventional and high-dimensional Bayesian optimization
\cite{lin2024smooth}.  Hypervolume-based methods instead select designs by how
much they are expected to expand the currently observed tradeoff region
\cite{daulton2020differentiable,ament2023unexpected}.  We study both families
because composite modeling should be tested as a representation choice rather
than as a feature of only one acquisition rule.

The key observation is that many scientific objectives are not measured
directly.  An experiment or simulator first returns a richer intermediate
response---a spectrum, trajectory, field, or collection of physical
contributions---and a known calculation reduces that response to the final
objectives.  Bayesian optimization of composite functions (BOCF) was introduced
for this setting with one scalar objective: model the intermediate response and
apply the known scalar reduction inside expected improvement
\cite{astudillo2019bayesian}.  Scientific design, however, often requires
several competing objectives computed from the same response.  Multi-objective
BOCF (MOBOCF) applies the same modeling idea before a multi-objective
acquisition or scalarization.  It changes the statistical representation, not
the experiment, and can use more of the information that a scientific
evaluation already provides.

The central criterion is equal score but different sensitivity.  Two distinct
intermediate responses may produce the same final objectives yet respond
differently when one component changes.  A final-objective model receives the
same response target for both states, although they imply different routes to
improvement.  This happens naturally with norms, products, oscillatory maps,
and maxima.  MOBOCF is
expected to help when such pairs occur and the intermediate response is easier
to learn than the reduced objective.  A decomposition alone is not enough: an
additive reduction with constant sensitivity provides a useful boundary case.

Prior studies have already demonstrated composite optimization with
multi-objective material and structural design, smooth Tchebycheff
scalarization, and high-dimensional optical design
\cite{coelho2025composite,pires2026composite,maddox2021physics}.  Building on
those methods, this paper is an application-focused investigation across four
established MOBO approaches and several scientific computation graphs.  We do
not propose a new acquisition function.  For each method, we compare an
objective-only model with a composite model while holding the initial designs,
acquisition rule, and expensive-evaluation budget fixed.  The study asks when
preserving intermediate scientific information improves tradeoff discovery in
low- and high-dimensional applications.

Our contributions are:
\begin{enumerate}
    \item We apply composite modeling within four existing MOBO approaches that
    cover hypervolume, scalarization, high-dimensional, and trust-region search.
    \item We assemble scientific and controlled synthetic applications around
    an equal-score/different-sensitivity criterion, with an additive boundary
    case.
    \item We isolate surrogate representation with paired initial designs and
    matched expensive-evaluation budgets, then evaluate final hypervolume,
    complete learning curves, and computational overhead across ten trials.
    \item We translate the comparisons into practical guidance about when a
    scientific workflow should retain and model its intermediate response.
\end{enumerate}

The experiments contain two controlled synthetic benchmarks and three
deterministic scientific simulators.  All have two objectives; their input
dimensions range from 4 to 600.  The low-dimensional study compares qLogEHVI
and smooth Tchebycheff search, while the high-dimensional study compares
spherical-linear Tchebycheff search and MORBO.  Across ten independent paired
trials, composite modeling improves all ten mean final hypervolumes.  The median
relative final gain is 6.98\%, with the largest benefits on strong instances of
the equal-score/different-sensitivity criterion.


\section{Related work}
\label{sec:related}

Bayesian optimization uses a probabilistic surrogate and an acquisition
function to allocate expensive evaluations \cite{shahriari2016taking,frazier2018tutorial}.
Modern MOBO methods either optimize an indicator such as dominated hypervolume
or repeatedly optimize scalar preferences.  Differentiable Monte Carlo
expected hypervolume improvement (qEHVI) \cite{daulton2020differentiable} and
its numerically stable logarithmic form \cite{ament2023unexpected} provide the
first route.  ParEGO \cite{knowles2006parego} and smooth Tchebycheff (STCH)
scalarization \cite{lin2024smooth} provide the second.  MORBO coordinates
multiple local trust regions to make Pareto discovery viable in high input
dimension \cite{daulton2022morbo}.  Recent work also shows that a geometric
warping followed by a linear kernel can be unexpectedly competitive in very
high dimensions \cite{doumont2026linear}.

BOCF instead models $h$ and induces a generally non-Gaussian posterior over
$\composemap(h)$.  Astudillo and Frazier introduced this idea for a single
scalar objective and developed expected improvement for composite functions
\cite{astudillo2019bayesian}; function-network models extend the same
single-objective principle to more general directed computations
\cite{astudillo2021networks}.  Structured-output models later showed that
composite BO can scale to responses with thousands of correlated entries
\cite{maddox2021outputs}.  Constrained composite BO has also been validated in
wet-lab polymer-particle synthesis, again for target-oriented scalar
optimization \cite{wang2025constrained}.

The extension to several competing objectives is newer but not absent.
Cardoso Coelho et al. study composite surrogates in multi-objective material
and structural design \cite{coelho2025composite}; Pires et al. combine
composite models with smooth Tchebycheff scalarization
\cite{pires2026composite}; and Maddox et al. combine a structured-output model
with MORBO for a 177-dimensional, two-objective diffractive-optics application
\cite{maddox2021physics}.  These works establish that MOBOCF can be implemented
with several surrogate and acquisition designs.  The present paper therefore
does not introduce composite qEHVI, STCH, or MORBO.  It contributes an
application-focused, matched evaluation: we apply the same representation
change across four existing MOBO regimes on scientific computation graphs
chosen by a common structural criterion.

\section{Problem setting}
\label{sec:problem}

This section defines the objective structure, observation assumptions, and
optimization goal used throughout the study.  We consider deterministic
minimization on a normalized design box
$\mathcal X=[0,1]^d$:
\begin{equation}
  \min_{x\in\mathcal X} f(x)
  =\min_{x\in\mathcal X}\composemap(h(x)),
  \qquad h:\mathcal X\rightarrow\R^K,
  \quad \composemap:\R^K\rightarrow\R^M.
  \label{eq:problem}
\end{equation}
Here $h(x)$ concatenates all intermediate responses needed by the $M$
objectives.  Components may be objective-specific, duplicated, or shared by
several objectives; the benchmark definitions state which case applies.  One
joint expensive evaluation observes all entries of $h(x)$ without noise, and
all final objectives are then computed by the deterministic, known, and
inexpensive map $\composemap$.  Every objective is minimized.  These assumptions
describe the controlled experiments in this paper; partial and noisy
intermediate measurements are left for future work.

\paragraph{Equal-score/different-sensitivity criterion.}
The computation graph satisfies our criterion when there are two attainable
intermediate states $u,v\in h(\mathcal X)$ such that
\[
u\ne v,\qquad \composemap(u)=\composemap(v),\qquad
J_{\composemap}(u)\ne J_{\composemap}(v),
\]
where $J_{\composemap}$ is the Jacobian of the known objective map.  In
numerical problems, equality may be interpreted within a fixed tolerance.  For
a nonsmooth map such as a maximum, the corresponding requirement is that the
active component or subgradient differs.  In plain terms, the same objective
value conceals different local responses to a component change.  This is
stronger than requiring the map to be merely many-to-one: an additive sum is
many-to-one but has constant sensitivity, so it does not satisfy the strong
criterion.  Observing $h$ and knowing $\composemap$ exactly are also required;
otherwise the composite model cannot use the distinction.

A design $x$ is Pareto dominated by $x'$ if
$f_m(x')\le f_m(x)$ for every $m$ and the inequality is strict for at least one
objective.  The Pareto set $\mathcal P_X$ contains designs that are not
dominated by any point in $\mathcal X$, and the Pareto front is
$\mathcal P_F=\{f(x):x\in\mathcal P_X\}$.  Given a minimization reference point
$r$, the quality of an observed nondominated set $P_n$ is its dominated
hypervolume
\begin{equation}
  \HV(P_n;r)=\lambda_M\!\left(
  \bigcup_{y\in P_n}[y_1,r_1]\times\cdots\times[y_M,r_M]
  \right),
  \label{eq:hv}
\end{equation}
where $\lambda_M$ is $M$-dimensional Lebesgue measure.  Hypervolume is the sole
performance metric used to rank the methods.

Every expensive evaluation reveals either $f(x)$ alone to a direct solver, or
$h(x)$ to a composite solver; in the latter case $f(x)$ is computed without a
second oracle call.  For a direct model we fit independent GPs
$\widehat f_1,\ldots,\widehat f_M$.  For a composite model we fit independent
GPs $\widehat h_1,\ldots,\widehat h_K$.  If
$H^{(s)}(x)\sim p(h(x)\mid\mathcal D_n)$ is a joint posterior draw, then
\begin{equation}
  F^{(s)}(x)=\composemap\!\left(H^{(s)}(x)\right)
  \label{eq:posterior-pushforward}
\end{equation}
is used inside the same acquisition function as the direct posterior draw.
Thus the expensive-evaluation interface, objective definition, and acquisition
criterion are unchanged; only the surrogate representation differs.  Our
independent component GPs do not exploit correlations among entries of $h$, so
the comparison tests a deliberately simple and broadly reusable form of
composition.

\section{Considered methods}
\label{sec:methods}

We evaluate four method families: qLogEHVI, smooth Tchebycheff (STCH),
spherical-linear STCH, and MORBO.  Every family has a matched pair that differs
only in its modeled response.  We use the names Direct--$X$ for the
objective-space version and Composite--$X$ for its composite counterpart.
Complete model and optimization settings appear in
Appendix~\ref{app:implementation}.

\subsection{qLogEHVI}
The direct method fits standardized independent GPs to $-f_m$ and chooses one
point at a time with qLogEHVI.  The composite method fits the same GP class to
$h_k$ and applies $-\composemap$ to posterior samples inside qLogEHVI.  Both use
the same observed Pareto partition and minimization reference point after the
sign conversion.

\subsection{Smooth Tchebycheff scalarization}
For an ideal point $z^\star$, positive preference vector $w$, and temperature
$\tau$, we minimize
\begin{equation}
  s_w(y)=\tau\log\sum_{m=1}^{M}
  \exp\!\left(\frac{w_m(y_m-z_m^\star)}{\tau}\right).
  \label{eq:stch}
\end{equation}
The direct method applies $-s_w$ to objective-posterior draws and uses
qLogExpectedImprovement.  The composite method applies
$-s_w\!\circ\composemap$ to component-posterior draws.  Five fixed weights
cover the two-objective simplex.  Each weight starts from the same initial
design, branches independently, and its new designs are charged to one common
evaluation budget.

\subsection{Spherical-linear STCH}
In high dimension we replace the usual stationary kernel with the
spherical-linear construction of Doumont et al. \cite{doumont2026linear}.
After centering and learned coordinate-wise scaling, $u\in\R^d$ is mapped to
the unit sphere by the inverse stereographic projection
\begin{equation}
  P(u)=\frac{(2u,\,\lVert u\rVert^2-1)}{1+\lVert u\rVert^2},
  \qquad
  k(u,u')=b_0+b_1P(u)^\top P(u'),
  \label{eq:spherical}
\end{equation}
with $(b_0,b_1)$ constrained to the simplex.  Direct and composite variants
then use the STCH procedure above.  This finite-feature model is intentionally
low capacity; composition asks it to learn smooth latent coordinates instead
of a highly nonlinear aggregate.

\subsection{MORBO}
MORBO maintains several local trust regions and coordinates their candidates
to improve the global Pareto set \cite{daulton2022morbo}.  We use the vendored
batched engine with five regions and joint batches of five.  The direct version
fits local GPs to the objectives.  The composite version fits the same local GP
class to intermediate components and composes Thompson posterior samples
before the engine's hypervolume-improvement scoring.  This is an adaptation of
MORBO's surrogate/output interface, not an untouched default implementation;
full settings are given in Appendix~\ref{app:implementation}.

\section{Numerical experiments}
\label{sec:experiments}

The experiments cover synthetic and scientific problems in low- and
high-dimensional design spaces.  Dominated hypervolume is the primary metric.
Every Composite--$X$ method is compared with its matched Direct--$X$ method
under the same expensive-evaluation budget and paired initial design.

\subsection{Test problems}

\paragraph{Benchmark inclusion criterion.}
We choose problems by the equal-score/different-sensitivity property defined in
Section~\ref{sec:problem}, not by their observed optimization result.  The same
expensive evaluation must also reveal $h$, and $\composemap$ must be known and
cheap.  Norms meet the criterion because equal lengths can have different
directions; products because equal products can have different factor
sensitivities; oscillatory maps because repeated values can occur at points
with different slopes; and maxima because the same maximum can be controlled
by different components.  Color formulation, DTLZ2, Langermann--Ackley, and
shallow lake are strong instances.  The truss is a planned boundary case: its
outer map is additive, so its sensitivity is constant.  This
gives a direct hypothesis: MOBOCF should help most on the four strong instances
and less on the additive boundary case, provided the intermediate components
can be learned within the budget.

Table~\ref{tab:benchmarks} summarizes the five two-objective minimization
problems.  The three low-dimensional cases compare qLogEHVI and STCH; the two
high-dimensional cases compare spherical-linear STCH and MORBO.  Color
formulation, truss design, and shallow-lake management are deterministic
scientific simulators; they are not tabular measured datasets.  DTLZ2 and
Langermann--Ackley are controlled synthetic compositions.  Detailed equations
and all modeled quantities appear in Appendix~\ref{app:benchmarks}.

\begin{table}[t]
\centering
\caption{Benchmark graphs and budgets.  $K$ is the number of independently
modeled intermediate outputs; all problems have $M=2$ minimized objectives.}
\label{tab:benchmarks}
\small
\setlength{\tabcolsep}{4.5pt}
\begin{tabular}{lcccll}
\toprule
Benchmark & Type & $d$ & $K$ & Initial/total & Compared families \\
\midrule
Color formulation & scientific & 6 & 6 & 5/45 & qLogEHVI, STCH \\
DTLZ2 & synthetic & 6 & 4 & 5/45 & qLogEHVI, STCH \\
Four-bar truss (RE21) & scientific & 4 & 8 & 5/45 & qLogEHVI, STCH \\
Langermann--Ackley & synthetic & 600 & 12 & 20/120 & spherical STCH, MORBO \\
Shallow lake & scientific & 100 & 21 & 20/120 & spherical STCH, MORBO \\
\bottomrule
\end{tabular}
\end{table}

\subsubsection{Synthetic problems}

\paragraph{DTLZ2.}
The six-dimensional DTLZ2 problem is the controlled low-dimensional synthetic
case \cite{deb2002scalable}.  Five variables determine a nonnegative distance
$g(x)$ from the Pareto set, while the remaining variable determines a
trigonometric position on the front.  The final objectives multiply $1+g(x)$
by cosine and sine position factors.  We expose four intermediates---separate
copies of $g$ and the two trigonometric factors---so the composite surrogates
learn distance and position before the known products are applied.  This
factorization cleanly isolates the proposed mechanism because the exact Pareto
front is known, yet direct objective models must learn the product geometry.

\paragraph{Latent Langermann--Ackley.}
This controlled high-dimensional problem maps 600 design coordinates through
six fixed dense orthonormal projections.  One objective applies an oscillatory
Langermann-style landscape to those latent coordinates, and the other applies
the nested exponential and cosine operations of Ackley.  Objective-specific
copies give $K=12$ independently modeled intermediates, deliberately avoiding
implicit information sharing.  Its nominal dimension is 600 but its effective
dimension is six, allowing us to ask whether component observations expose a
learnable low-dimensional response that is hidden by two highly nonlinear
outer maps.

\subsubsection{Scientific problems}

\paragraph{Multi-illuminant color formulation.}
This six-variable scientific benchmark represents the concentrations of six
hypothetical dyes in an opaque coating.  The simulator first converts the dye
mixture into a 31-wavelength reflectance spectrum with the Kubelka--Munk model,
then converts that spectrum into CIE Lab coordinates under a daylight
approximation and Illuminant A.  These six illuminant-specific coordinates are
the observed intermediate response.  The two minimized objectives are a
daylight color-error norm and a metameric-mismatch norm measuring how the
sample's error changes between illuminants.  A direct model therefore sees
only two scalar distances, whereas the composite model retains the direction
of the color error and the illumination under which it occurs.  This is
representative of materials workflows in which spectroscopy or
characterization produces a structured response before a small collection of
application scores is computed \cite{chaoucha2022color,kubelka1931optik}.

\paragraph{Four-bar truss.}
The four design variables in RE21 are structural member cross-sectional areas
subject to allowable-stress bounds \cite{tanabe2020realworld}.  Eight
intermediates are the four member-volume contributions and four signed
compliance contributions.  Their known outer map sums and normalizes these
terms into total material volume and loaded-joint displacement.  Unlike the
other reductions, this map is essentially additive and monotone.  The problem
therefore serves as a scientific boundary case: its component values are
physically interpretable, but the final objectives themselves remain smooth
and do not strongly discard component geometry.

\paragraph{Shallow-lake management.}
The 100 design variables specify annual phosphorus-release policies in a
century-long dynamical simulation \cite{carpenter1999lake,ward2015tipping}.
Across 100 fixed inflow scenarios, the simulator returns 20 normalized
five-year phosphorus peaks and one normalized discounted-utility value.  The
objectives minimize the worst block peak and one minus utility.  Taking a
maximum creates a nonsmooth envelope whenever the controlling time block
changes and erases which part of the trajectory caused the objective.  The
composite model instead observes the time-resolved block summaries before
applying this maximum, mirroring scientific control problems in which an
entire trajectory is available even though optimization is scored by a worst
event and an aggregate benefit.

\subsection{Evaluation metric and experimental settings}

We run 10 independent paired trials with seeds
$0,10007,\ldots,90063$.  A scrambled Sobol design supplies 5 initial points in
the low-dimensional suite and 20 in the high-dimensional suite.  The budgets
are 45 and 120 total expensive design evaluations, respectively; hence the
adaptive phases contain 40 and 100 evaluations.  For each method family, the
direct and composite variants have bit-identical initial designs.  Initial
designs need not match across different families.

For STCH the five weights are
$(0.05,0.95)$, $(0.275,0.725)$, $(0.5,0.5)$,
$(0.725,0.275)$, and $(0.95,0.05)$, with $\tau=0.05$.  Each low-dimensional
branch receives 8 adaptive evaluations; each high-dimensional branch receives
20.  The plotted trace interleaves branches in round-robin weight order after
the shared initial design.  MORBO performs 20 joint batches of size 5 after its
20 initial points.  Low-dimensional acquisition optimization uses 128 raw
samples and 8 restarts; high-dimensional STCH uses 256 and 10.  All acquisition
optimizers use at most 200 iterations.

We report dominated hypervolume, not hypervolume regret.  Objectives are scaled
to a nominal $[0,1]$ range.  The reference point is $(1.1,1.1)$ for DTLZ2 and
Langermann--Ackley and $(1.05,1.05)$ otherwise.  No known Pareto front is used
to compute the metric; the analytic DTLZ2 ceiling is shown only for
interpretation.  Because reference points and objective meanings differ,
absolute hypervolumes are compared only within a benchmark.  We report the
mean final hypervolume and its standard error over trials, the number of paired
trials in which composite is larger, and the relative change
$(\overline{\HV}_{\cf}-\overline{\HV}_{\direct})/
\overline{\HV}_{\direct}$.  We also average the hypervolume over every
post-initial evaluation to obtain a time-averaged learning-curve summary.

\section{Results and discussion}
\label{sec:results}

Composite modeling increases mean final hypervolume in every comparison in
Table~\ref{tab:results}; Figures~\ref{fig:low-learning-curves}
and~\ref{fig:high-learning-curves} show the complete learning curves.  It wins
89 of 100 individual paired trials.  Across
the ten comparisons, the median relative final gain is 6.98\%, and the median
post-initial time-averaged gain is 4.98\%.  The scale of benefit is strongly
problem- and method-dependent: the largest gain is 69.91\% for STCH on DTLZ2,
whereas the additive four-bar graph yields less than 1\% at the final budget.

\begin{table}[t]
\centering
\caption{Final dominated hypervolume over 10 paired trials (mean $\pm$ standard
error).  Gain is computed from the displayed methods' unrounded means.}
\label{tab:results}
\footnotesize
\setlength{\tabcolsep}{3.8pt}
\begin{tabular}{llrrr}
\toprule
Benchmark & Method family & Direct & Composite & Gain \\
\midrule
Color formulation & qLogEHVI
  & $0.8502\!\pm\!0.0110$ & $0.9672\!\pm\!0.0027$ & $+13.76\%$ \\
 & STCH
  & $0.8045\!\pm\!0.0124$ & $0.8747\!\pm\!0.0119$ & $+8.73\%$ \\
DTLZ2 & qLogEHVI
  & $0.2928\!\pm\!0.0311$ & $0.4053\!\pm\!0.0008$ & $+38.43\%$ \\
 & STCH
  & $0.0839\!\pm\!0.0185$ & $0.1425\!\pm\!0.0193$ & $+69.91\%$ \\
Four-bar truss & qLogEHVI
  & $0.8503\!\pm\!0.0006$ & $0.8573\!\pm\!0.0001$ & $+0.82\%$ \\
 & STCH
  & $0.8095\!\pm\!0.0030$ & $0.8151\!\pm\!0.0014$ & $+0.69\%$ \\
\midrule
Langermann--Ackley & spherical STCH
  & $0.6650\!\pm\!0.0098$ & $0.6888\!\pm\!0.0072$ & $+3.59\%$ \\
 & MORBO
  & $0.6368\!\pm\!0.0139$ & $0.6705\!\pm\!0.0154$ & $+5.29\%$ \\
Shallow lake & spherical STCH
  & $0.1071\!\pm\!0.0009$ & $0.1164\!\pm\!0.0008$ & $+8.66\%$ \\
 & MORBO
  & $0.0954\!\pm\!0.0010$ & $0.0993\!\pm\!0.0013$ & $+4.08\%$ \\
\bottomrule
\end{tabular}
\end{table}

\subsection{Low-dimensional problems}

\begin{figure}[!ht]
\centering
\begin{minipage}[t]{0.49\linewidth}
  \centering
  \resultplot[width=\linewidth]{color_formulation_graph.png}
  \par\vspace{0.4em}
  {\small (a) Color formulation\par}
\end{minipage}\hfill
\begin{minipage}[t]{0.49\linewidth}
  \centering
  \resultplot[width=\linewidth]{dtlz2_graph.png}
  \par\vspace{0.4em}
  {\small (b) DTLZ2\par}
\end{minipage}\par\vspace{1.25em}
\begin{minipage}[t]{0.58\linewidth}
  \centering
  \resultplot[width=\linewidth]{four_bar_truss_graph.png}
  \par\vspace{0.4em}
  {\small (c) Four-bar truss RE21\par}
\end{minipage}
\caption{Low-dimensional dominated-hypervolume learning curves.  Lines are
means and bands are one standard error over 10 trials; the vertical line ends
the 5-point initial design.  Each method uses 45 total evaluations.}
\label{fig:low-learning-curves}
\end{figure}

DTLZ2 provides the clearest low-dimensional instance of the inclusion
criterion.  Its direct objectives
multiply a distance term by trigonometric position terms; the composite model
instead observes those smooth factors.  With reference point $(1.1,1.1)$, the
exact attainable hypervolume is
$1.1^2-\pi/4=0.42460$.  Composite qLogEHVI reaches 0.40535 on average
(95.46\% of the ceiling), compared with 0.29281 (68.96\%) for direct qLogEHVI.
Color formulation also exhibits a substantial advantage.  Its two scores are
norms of six illuminant-specific Lab coordinates, so the outer reductions erase
both direction and illuminant-specific error; preserving those coordinates
raises final hypervolume by 13.76\% under qLogEHVI and 8.73\% under STCH.

The four-bar truss has the same formal composite graph but a nearly additive,
monotone reduction.  Its direct objectives are already smooth and easy to fit,
so final gains are only 0.82\% and 0.69\%.  Composition improves early search
more clearly: its post-initial time-averaged gains are 2.91\% for qLogEHVI and
5.16\% for STCH.  This contrast supports a more specific hypothesis than
``more outputs are better'': composition is most useful when the known outer
same final value can correspond to factor states with different sensitivities.
Thus the composite advantage appears both during the adaptive phase and at the
final budget, and it occurs under both hypervolume and scalarization search.

\subsection{High-dimensional problems}

\begin{figure}[!ht]
\centering
\begin{minipage}[t]{0.49\linewidth}
  \centering
  \resultplot[width=\linewidth]{langermann_ackley_graph.png}
  \par\vspace{0.4em}
  {\small (a) Langermann--Ackley\par}
\end{minipage}\hfill
\begin{minipage}[t]{0.49\linewidth}
  \centering
  \resultplot[width=\linewidth]{shallow_lake_graph.png}
  \par\vspace{0.4em}
  {\small (b) Shallow lake\par}
\end{minipage}
\caption{High-dimensional dominated-hypervolume learning curves.  Lines are
means and bands are one standard error over 10 trials; the vertical line ends
the 20-point initial design.  Each method uses 120 total evaluations.}
\label{fig:high-learning-curves}
\end{figure}

In 600 dimensions, Langermann--Ackley depends on six dense orthonormal
projections.  Direct surrogates must represent an oscillatory Langermann map
and the nested exponential Ackley map over the ambient box.  Composite
spherical-linear GPs see the projections themselves, increasing final
hypervolume by 3.59\%; component-level local GPs raise MORBO by 5.29\%.
The gains are meaningful but smaller than DTLZ2 because the twelve independent
component GPs duplicate the same six latent coordinates and do not share
information across objectives.

The shallow-lake outer objective takes a maximum over 20 time blocks.  The
identity of the controlling block changes with the release policy, creating a
nonsmooth envelope and discarding temporal location.  Modeling block peaks
before the maximum improves spherical STCH by 8.66\% and MORBO by 4.08\%.
This is the clearest scientific high-dimensional example of equal score but
different sensitivity: the same worst value can be controlled by a different
part of the trajectory.

All four high-dimensional method--benchmark pairs improve at the final budget,
with a median relative gain of 4.69\%.  Their post-initial time-averaged gains
range from 2.59\% to 6.67\%, so the benefit also appears before the final
evaluation.  These results do not show that input dimension alone increases
the value of composition.  Instead, spherical-linear modeling and MORBO handle
the large design space, while the composite representation determines whether
their models retain useful response information.

\subsection{When does composite modeling help?}
\label{sec:discussion}

The results follow the equal-score/different-sensitivity criterion.  The color
norms assign the same distance to errors pointing in different directions;
the DTLZ2 products can match while their factors have different derivatives;
the Langermann--Ackley maps repeat values with different slopes; and the
shallow-lake maximum can switch its controlling time block.  In each case, the
composite model distinguishes states that an objective-only model merges.  The
number of components must still be small or structured enough to model within
the available budget.  This explains why the gains vary: DTLZ2 and color are
strong instances, while the additive truss has constant outer sensitivity and
yields less than one percent at the final budget.  More outputs
alone are not the point; equal-score states must imply meaningfully different
local behavior.

This interpretation is consistent across, but not identical for, the four
optimization regimes.  Direct and composite variants use the same acquisition
principle, initialization, and expensive-evaluation budget within each pair,
so the paired difference isolates which response is modeled.  Improvements
under qLogEHVI, conventional STCH, spherical STCH, and MORBO indicate that the
effect is not confined to one scalarization or one global optimizer.  At the
same time, the gains differ between method families on the same problem.  The
representation and the acquisition geometry therefore interact: a more
informative posterior can help several acquisition rules, but each rule uses
that posterior differently.  Establishing causal contributions from
particular map properties would require a dedicated ablation that varies one
property at a time, so we treat the product, norm, oscillation, and maximum
explanations as mechanisms supported by the controlled contrasts rather than
as separately identified effects.

\paragraph{Implications for scientific design.}
The practical opportunity is largest when an experiment or simulator already
returns a rich characterization and the reported objectives are deterministic
summaries of it.  A spectrometer returns a curve before software computes a
color error; a dynamical model returns a trajectory before a worst-time score;
and a mechanics calculation returns member-level quantities before total
volume or displacement.  Discarding these responses forces a direct GP to
reconstruct the geometry of the summary from one or two numbers per design.
Composite MOBO keeps the same number of synthesized samples or simulator calls
while retaining the response that those calls already produced.  This makes
the approach especially relevant to the AI-guided-design part of a
self-driving laboratory: it changes the statistical interface between
characterization and decision-making without requiring a new instrument or an
additional physical experiment
\cite{macleod2022selfdriving,gopakumar2018multiobjective}.

The results also suggest a practical screening rule.  Composite modeling is a
strong candidate when the intermediate response is observable at negligible
additional experimental cost, its transformation into objectives is known and
stable, and that transformation contains a norm, extremum, product, threshold,
or another reduction that loses identifiable structure.  A direct model may
be preferable when the final objectives are already smooth and low-complexity,
when the component vector is extremely large relative to the budget, or when
the supposed components are noisy and only weakly related to Pareto
improvement.  The four-bar result illustrates the first case, while the
duplicated Langermann coordinates and 21-output lake model illustrate the
statistical and computational price of the second.  In practice, a short
paired pilot can compare posterior calibration and early hypervolume progress
before committing the full experimental campaign.

\paragraph{Sample efficiency and computational cost.}
Composition adds no physical evaluations in these comparisons, but it does not
make optimization free.  The median post-initial hypervolume gain of 4.98\%
shows that the advantage is present across the trajectory rather than only at
the final design.  Against this, composite runs take a median 3.45 times as
long, with the largest overhead for the 21-output shallow-lake spherical
model.  Whether that exchange is favorable depends on the scientific cost
model.  Several extra hours of surrogate fitting may be negligible when a
synthesis-and-characterization cycle takes a day and consumes scarce material;
the same overhead may be unacceptable for a cheap analytic simulator.  Better
multi-output kernels, shared latent models, and selective component modeling
could reduce this cost because the present implementation fits independent GPs
even when components are correlated or duplicated.

\subsection{Limitations and next steps}
The suite is intentionally scoped to the equal-score/different-sensitivity
criterion above.  It therefore does not establish universal or average
superiority over arbitrary black-box objectives; it evaluates MOBOCF where
there is a clear reason for preserving the intermediate response.  The
reported oracles are deterministic simulators rather than new physical
experiments, and the color and shallow-lake cases are simplified slices of
richer scientific domains.  The high-dimensional benchmarks also have
exploitable structure: Langermann--Ackley has six-dimensional effective
variation, while the lake objective is organized by time.  Hyperparameters and
acquisition settings are matched but not exhaustively tuned, runtime
measurements are cluster wall times rather than hardware-normalized complexity
experiments, and every problem has two objectives and a box domain.  The study
does not yet address measurement noise, unknown or uncertain outer maps,
missing intermediate observations, feasibility constraints, mixture-simplex
domains, or more than two target properties.

A broader study should compare direct, independent-component, and correlated
multi-output surrogates across a wider range of transformations, including
maps that preserve most intermediate information.  It should also allow the optimizer to choose
adaptively between direct and composite representations as evidence
accumulates, and should test partial or noisy component observations and
learned outer maps.  Most importantly, a prospective wet-lab campaign should
measure whether the hypervolume benefit survives calibration drift and
experimental variability.  Hyperspectral nanoparticle synthesis
\cite{low2024constrained} and composition or process optimization in which
spectroscopy is already part of the automated loop are natural candidates:
they directly test whether retaining characterization outputs accelerates the
discovery of Pareto-improving materials.

\section{Conclusion}
\label{sec:conclusion}

This application study asked whether established multi-objective Bayesian
optimization methods should model only the final scientific objectives or
preserve observable intermediate responses and apply the known objective map
afterward.  Across four strong instances of the
equal-score/different-sensitivity criterion, one additive boundary case, and
four MOBO regimes, the composite
variant improved all ten mean final hypervolumes under matched expensive
evaluation budgets.  The median relative final gain was 6.98\%, and the median
post-initial time-averaged gain was 4.98\%.  The result spans low-dimensional
global qLogEHVI and STCH, high-dimensional spherical STCH, and a batched MORBO
wrapper, showing that the established BOCF representation can benefit both
hypervolume- and scalarization-based scientific search.

The comparison also identifies the scope of that conclusion.  The largest
benefits arise when equal objective values correspond to different local
sensitivities: equal norms can point in different directions, equal products
can have different factors, repeated oscillatory values can have different
slopes, and equal maxima can be controlled by different time blocks.  A
physically meaningful decomposition is not automatically useful; when the
outer map has constant sensitivity, as in the four-bar truss, direct objective
models can already capture most of the relevant geometry.  Composition should
therefore be chosen from this criterion and the available budget, not treated
as a default replacement for direct MOBO.

For automated scientific and materials design, the central implication is
simple: characterization data should not be collapsed prematurely when the
full response is already available.  Composite MOBO can extract more learning
signal from each physical evaluation without asking the laboratory to perform
more experiments, although the additional surrogate fitting must be justified
by the cost of those experiments.  The next step is to move beyond
deterministic two-objective simulators toward correlated multi-output models,
partial and noisy responses, uncertain objective maps, adaptive selection
between direct and composite surrogates, and prospective validation in an
automated materials campaign.  Under those conditions, composition is best
viewed not as a new experimental instrument, but as a way to stop discarding
information produced by instruments and simulators already inside the
self-driving laboratory.

\clearpage
\bibliographystyle{plain}
\begin{thebibliography}{99}

\bibitem{ament2023unexpected}
S.~Ament, S.~Daulton, D.~Eriksson, M.~Balandat, and E.~Bakshy.
Unexpected improvements to expected improvement for Bayesian optimization.
In \emph{Advances in Neural Information Processing Systems}, 36, 2023.

\bibitem{astudillo2019bayesian}
R.~Astudillo and P.~I. Frazier.
Bayesian optimization of composite functions.
In \emph{Proceedings of the 36th International Conference on Machine Learning},
PMLR 97:354--363, 2019.

\bibitem{astudillo2021networks}
R.~Astudillo and P.~I. Frazier.
Bayesian optimization of function networks.
In \emph{Advances in Neural Information Processing Systems}, 34, 2021.

\bibitem{balandat2020botorch}
M.~Balandat, B.~Karrer, D.~R. Jiang, S.~Daulton, B.~Letham, A.~G. Wilson, and E.~Bakshy.
BoTorch: A framework for efficient Monte-Carlo Bayesian optimization.
In \emph{Advances in Neural Information Processing Systems}, 33, 2020.

\bibitem{carpenter1999lake}
S.~R. Carpenter, D.~Ludwig, and W.~A. Brock.
Management of eutrophication for lakes subject to potentially irreversible change.
\emph{Ecological Applications}, 9(3):751--771, 1999.
doi:10.1890/1051-0761(1999)009[0751:MOEFLS]2.0.CO;2.

\bibitem{chaoucha2022color}
S.~Chaoucha, A.~Moussa, and N.~Ladhari.
Color formulation of cotton fabrics using multi-objective ant colony optimization.
\emph{Journal of Natural Fibers}, 19(17):15459--15474, 2022.
doi:10.1080/15440478.2022.2128145.

\bibitem{coelho2025composite}
R.~P. Cardoso Coelho, A.~F. Carvalho Alves, T.~M. Nogueira Pires, and
F.~M. Andrade Pires.
A composite Bayesian optimisation framework for material and structural design.
\emph{Computer Methods in Applied Mechanics and Engineering}, 434:117516, 2025.
doi:10.1016/j.cma.2024.117516.

\bibitem{daulton2020differentiable}
S.~Daulton, M.~Balandat, and E.~Bakshy.
Differentiable expected hypervolume improvement for parallel multi-objective
Bayesian optimization.
In \emph{Advances in Neural Information Processing Systems}, 33, 2020.

\bibitem{daulton2022morbo}
S.~Daulton, D.~Eriksson, M.~Balandat, and E.~Bakshy.
Multi-objective Bayesian optimization over high-dimensional search spaces.
In \emph{Proceedings of the 38th Conference on Uncertainty in Artificial
Intelligence}, PMLR 180:507--517, 2022.

\bibitem{deb2002scalable}
K.~Deb, L.~Thiele, M.~Laumanns, and E.~Zitzler.
Scalable multi-objective optimization test problems.
In \emph{Proceedings of the 2002 Congress on Evolutionary Computation},
pages 825--830, 2002.

\bibitem{doumont2026linear}
C.~Doumont, D.~Fan, N.~Maus, J.~R. Gardner, H.~Moss, and G.~Pleiss.
We still don't understand high-dimensional Bayesian optimization.
In \emph{Proceedings of the 29th International Conference on Artificial
Intelligence and Statistics}, 2026.

\bibitem{frazier2018tutorial}
P.~I. Frazier.
A tutorial on Bayesian optimization.
\emph{arXiv:1807.02811}, 2018.

\bibitem{gopakumar2018multiobjective}
A.~M. Gopakumar, P.~V. Balachandran, D.~Xue, J.~E. Gubernatis, and T.~Lookman.
Multi-objective optimization for materials discovery via adaptive design.
\emph{Scientific Reports}, 8:3738, 2018.
doi:10.1038/s41598-018-21936-3.

\bibitem{knowles2006parego}
J.~Knowles.
ParEGO: A hybrid algorithm with on-line landscape approximation for expensive
multiobjective optimization problems.
\emph{IEEE Transactions on Evolutionary Computation}, 10(1):50--66, 2006.

\bibitem{kubelka1931optik}
P.~Kubelka and F.~Munk.
Ein Beitrag zur Optik der Farbanstriche.
\emph{Zeitschrift f\"ur Technische Physik}, 12:593--601, 1931.

\bibitem{lin2024smooth}
X.~Lin, Z.~Yang, X.~Zhang, and Q.~Zhang.
Smooth Tchebycheff scalarization for multi-objective optimization.
In \emph{Proceedings of the 41st International Conference on Machine Learning},
2024.

\bibitem{lookman2019active}
T.~Lookman, P.~V. Balachandran, D.~Xue, and R.~Yuan.
Active learning in materials science with emphasis on adaptive sampling using
uncertainties for targeted design.
\emph{npj Computational Materials}, 5:21, 2019.
doi:10.1038/s41524-019-0153-8.

\bibitem{low2024constrained}
A.~K. Low et al.
Evolution-guided Bayesian optimization for constrained multi-objective
optimization in self-driving labs.
\emph{npj Computational Materials}, 10:104, 2024.
doi:10.1038/s41524-024-01274-x.

\bibitem{macleod2022selfdriving}
B.~P. MacLeod et al.
A self-driving laboratory advances the Pareto front for material properties.
\emph{Nature Communications}, 13:995, 2022.
doi:10.1038/s41467-022-28580-6.

\bibitem{maddox2021outputs}
W.~J. Maddox, M.~Balandat, A.~G. Wilson, and E.~Bakshy.
Bayesian optimization with high-dimensional outputs.
In \emph{Advances in Neural Information Processing Systems}, 34, 2021.

\bibitem{maddox2021physics}
W.~J. Maddox, Q.~Feng, and M.~Balandat.
Optimizing high-dimensional physics simulations via composite Bayesian optimization.
In \emph{Fourth Workshop on Machine Learning and the Physical Sciences}, 2021.
arXiv:2111.14911.

\bibitem{pires2026composite}
T.~M. Nogueira Pires, R.~P. Cardoso Coelho, and F.~M. Andrade Pires.
Composite Bayesian optimisation for multi-objective problems with smooth
Tchebycheff scalarisation.
\emph{Structural and Multidisciplinary Optimization}, 69:34, 2026.
doi:10.1007/s00158-025-04229-y.

\bibitem{shahriari2016taking}
B.~Shahriari, K.~Swersky, Z.~Wang, R.~P. Adams, and N.~de~Freitas.
Taking the human out of the loop: A review of Bayesian optimization.
\emph{Proceedings of the IEEE}, 104(1):148--175, 2016.

\bibitem{stach2021autonomous}
E.~Stach et al.
Autonomous experimentation systems for materials development: A community perspective.
\emph{Matter}, 4(9):2702--2726, 2021.
doi:10.1016/j.matt.2021.06.036.

\bibitem{tanabe2020realworld}
R.~Tanabe and H.~Ishibuchi.
An easy-to-use real-world multi-objective optimization problem suite.
\emph{Applied Soft Computing}, 89:106078, 2020.
doi:10.1016/j.asoc.2020.106078.

\bibitem{wang2025constrained}
F.~Wang, M.~Parhizkar, A.~Harker, and M.~Edirisinghe.
Constrained composite Bayesian optimization for rational synthesis of polymeric particles.
\emph{Digital Discovery}, 4:3066--3077, 2025.
doi:10.1039/D5DD00243E.

\bibitem{ward2015tipping}
V.~L. Ward, R.~Singh, P.~M. Reed, and K.~Keller.
Confronting tipping points: Can multi-objective evolutionary algorithms discover
pollution control tradeoffs given environmental thresholds?
\emph{Environmental Modelling \& Software}, 73:27--43, 2015.
doi:10.1016/j.envsoft.2015.07.020.

\bibitem{wyman2013simple}
C.~Wyman, P.-P. Sloan, and P.~Shirley.
Simple analytic approximations to the CIE XYZ color matching functions.
\emph{Journal of Computer Graphics Techniques}, 2(2):1--11, 2013.

\end{thebibliography}

\appendix

\section{Detailed method descriptions}
\label{app:implementation}

\subsection{Shared direct/composite mechanics}

All calculations use double precision.  Inputs are normalized to $[0,1]^d$.
Each scalar output has its own \texttt{SingleTaskGP} and standardization
transform; collections are represented by a \texttt{ModelListGP}.  Unless
specified otherwise, BoTorch's default single-task GP covariance and inferred
Gaussian likelihood are used \cite{balandat2020botorch}.  Direct models have
$M=2$ outputs; composite models have the benchmark-specific $K$ outputs in
Table~\ref{tab:benchmarks}.  We deliberately do not use multi-task covariance,
latent sharing, or a component bottleneck.

At every proposed composite point, the oracle is called once for $h(x)$ and
$f(x)$ is reconstructed as $\composemap(h(x))$.  The implementation checks the
composition on domain validation and the archived results were re-evaluated
against the current benchmark code: objective and component reconstructions
agree to numerical precision (maximum discrepancies on the order of
$10^{-13}$).  Evaluation caches are cleared between methods, and every returned
history is asserted to contain exactly 45 or 120 designs.

The qLogEHVI run is sequential ($q=1$).  Acquisition optimization uses box
bounds, multistart gradient optimization, and a fresh Sobol fallback if a
nonlinear composition produces non-finite gradients.  STCH branches share only
their initial observations; after branching, the GP for each weight is updated
using that weight's own new points.  The archive combines the five branches in
round-robin order solely to define an order-neutral hypervolume trace.

\subsection{Spherical-linear surrogate}

For each output, the implementation centers $x$ at $(1/2,\ldots,1/2)$, divides
coordinates by learned ARD scales, and applies a learned global radial scale
before Eq.~\eqref{eq:spherical}.  A softmax constrains $b_0+b_1=1$ and
$b_0,b_1>0$.  The GP has a constant mean, a learned Gaussian likelihood with a
log-normal noise prior, and standardized outputs.  This is a finite
$O(d)$-feature linear model; posterior inference is $O(nd^2)$ and therefore
linear in the number of observations $n$ for fixed $d$, rather than an RBF GP
claiming to resolve a 600-dimensional surface from 120 points.

\subsection{MORBO wrapper}

The high-dimensional runner calls \texttt{batched\_morbo} and
\texttt{composite\_batched\_morbo}, which wrap the vendored coordinated engine
under \texttt{morbo/}; it does not call the separate sequential prototype.
There are five trust regions, each initialized with side length $0.8$, minimum
$0.01$, and maximum $1.6$.  Candidate pools contain 512 raw points.  Each local
output GP uses a Mat\'ern-$5/2$ ARD kernel with length scales constrained to
$[0.05,4]$, output scale, and a Gaussian likelihood with a Gamma noise prior.
The engine uses its lower-tail winsorization and coordinated Thompson/HVI
selection.  Minimum local training size is 19 because the initial design has
20 points.  The success streak is set to 10,000, effectively disabling region
expansion within this budget; failure streaks are
$\max(\lfloor d/3\rfloor,10)$, i.e., 200 for Langermann--Ackley and 33 for the
lake.  Twenty joint batches of five bring each history from 20 to 120 designs.

These details matter when interpreting ``MORBO'': this is a budget-adapted
wrapper around the published engine, not a claim that all published default
settings were preserved.  Direct and composite variants use identical engine
settings, which is sufficient for the representation comparison made here.

\section{Benchmark definitions}
\label{app:benchmarks}

This section specifies the exact deterministic oracle used in the experiments.
All vectors $x$ below are normalized to $[0,1]^d$ and all final objectives are
minimized.

\subsection{Multi-illuminant color formulation: $d=6$, $K=6$}

This analytical simulator represents six hypothetical dye loadings in an
opaque textile coating.  It is motivated by multi-objective color formulation
under multiple illuminants \cite{chaoucha2022color}, but is not a measured
dataset and does not reproduce that paper's objective equations.  Wavelengths
are $\lambda\in\{400,410,\ldots,700\}$ nm.  Define
\begin{equation}
 B(\lambda;\mu,\sigma)=
 \exp\!\left[-\tfrac12\left(\frac{\lambda-\mu}{\sigma}\right)^2\right].
\end{equation}
The substrate absorption-to-scattering ratio is
$q_0(\lambda)=0.012+0.008B(\lambda;430,75)$.  The six fixed dye increments are
\begin{equation}
\begin{aligned}
q_1 &= 1.75B(610,34)+0.22B(445,52),\\
q_2 &= 1.55B(535,31)+0.18B(650,42),\\
q_3 &= 1.65B(448,28)+0.16B(585,65),\\
q_4 &= 1.20B(495,48)+0.38B(632,27),\\
q_5 &= 1.35B(575,41)+0.24B(420,24),\\
q_6 &= 0.55+0.22B(510,115),
\end{aligned}
\end{equation}
where wavelength arguments are suppressed.  Concentrations are
$c_i=0.85x_i$ and $q=q_0+\sum_i c_iq_i$.  The opaque Kubelka--Munk inverse
\cite{kubelka1931optik} gives reflectance
\begin{equation}
  R(\lambda)=1+q(\lambda)-\sqrt{q(\lambda)^2+2q(\lambda)}.
  \label{eq:km}
\end{equation}

Analytic CIE 1931 matching curves from Wyman et al.
\cite{wyman2013simple} convert $R$ to XYZ under blackbody spectra at 6504 K
(daylight approximation, $D$) and 2856 K (Illuminant A).  With the illuminant's
perfect-reflector white point, XYZ is converted to CIE Lab using
\begin{equation}
 L^*=116\phi(Y/Y_n)-16,\quad
 a^*=500[\phi(X/X_n)-\phi(Y/Y_n)],\quad
 b^*=200[\phi(Y/Y_n)-\phi(Z/Z_n)],
\end{equation}
where $\delta=6/29$ and
$\phi(t)=t^{1/3}$ for $t>\delta^3$, otherwise
$t/(3\delta^2)+4/29$.  The six intermediate outputs are
\begin{equation}
 h(x)=\big(L_D^*,a_D^*,b_D^*,L_A^*,a_A^*,b_A^*\big).
\end{equation}

The unattainable target pigment has
\begin{equation}
q_\star=q_0+0.92B(482,19)+0.74B(622,25)+0.18B(550,72),
\end{equation}
and is converted to target vectors $t_D,t_A$ by the same pipeline.  Let
$e_D=h_{1:3}-t_D$ and $e_A=h_{4:6}-t_A$.  The objectives are
\begin{equation}
  f_1=\frac{\lVert e_D\rVert_2}{60},
  \qquad
  f_2=\frac{\lVert e_A-e_D\rVert_2}{18}.
  \label{eq:color-objectives}
\end{equation}
Thus $f_1$ is daylight color error and $f_2$ is a metameric-mismatch proxy: it
penalizes how the sample-minus-target Lab error changes between illuminants.
The nonlinear norms erase the direction of the six-dimensional color errors.
The hypervolume reference is $(1.05,1.05)$.

\subsection{DTLZ2: $d=6$, $K=4$}

The two-objective DTLZ2 problem \cite{deb2002scalable} is
\begin{equation}
 g(x)=\sum_{j=2}^{6}(x_j-0.5)^2,\qquad
 \theta(x)=\frac{\pi x_1}{2},
\end{equation}
\begin{equation}
 h(x)=\big(g,\cos\theta,g,\sin\theta\big),\qquad
 f_1=(1+h_1)h_2,\quad f_2=(1+h_3)h_4.
 \label{eq:dtlz}
\end{equation}
The repeated $g$ deliberately gives each objective an independently modeled
component group, avoiding implicit multi-task transfer.  The Pareto set has
$x_j=0.5$ for $j=2,\ldots,6$ and the front satisfies
$f_1^2+f_2^2=1$, $f_1,f_2\ge0$.  For reference point $(1.1,1.1)$ the exact
maximum hypervolume is $1.1^2-\pi/4$.

\subsection{Four-bar truss RE21: $d=4$, $K=8$}

This transparent structural-design problem follows the RE21 formulation in
the real-world multi-objective suite \cite{tanabe2020realworld}.  The four
cross-sectional areas are
\begin{equation}
 a=a^{\min}+x\odot(a^{\max}-a^{\min}),\quad
 a^{\min}=(1,\sqrt2,\sqrt2,1),\quad a^{\max}=(3,3,3,3)
\end{equation}
in cm$^2$.  These bounds encode the allowable-stress restriction for load
$F=10$ kN and allowable stress $10$ kN/cm$^2$.  Member length is $L=200$ cm
and elastic modulus is $E=2\times10^5$ kN/cm$^2$.  The component vector is
\begin{equation}
h=\left(2a_1,\sqrt2a_2,\sqrt{a_3},a_4,
        \frac2{a_1},\frac{2\sqrt2}{a_2},
        -\frac{2\sqrt2}{a_3},\frac2{a_4}\right).
\end{equation}
Raw volume and loaded-joint displacement are
\begin{equation}
 V=L\sum_{k=1}^4h_k,
 \qquad
 D=\frac{FL}{E}\sum_{k=5}^8h_k.
\end{equation}
They are affinely normalized using exact extrema over the box:
\begin{align}
V_{\min}&=1237.8414,& V_{\max}&=2994.9383,\\
D_{\min}&=0.00276142,& D_{\max}&=0.05057191,\\
f_1&=(V-V_{\min})/(V_{\max}-V_{\min}),&
f_2&=(D-D_{\min})/(D_{\max}-D_{\min}).
\end{align}
The reference is $(1.05,1.05)$.  Unlike the color and lake reductions, the
outer map is essentially a signed sum plus affine scaling; this explains why
the two objective GPs already capture most useful regularity.

\subsection{Latent Langermann--Ackley: $d=600$, $K=12$}

Let $A\in\R^{6\times600}$ have fixed orthonormal rows generated once with seed
18,271, and define the dense latent vector
\begin{equation}
 z=8A(x-\tfrac12\mathbf1).
\end{equation}
The component response contains objective-specific copies,
$h=(z,z)\in\R^{12}$.  This prevents the independent GP implementation from
sharing the six latent observations across objectives.  For centers
\begin{equation}
C=\begin{bmatrix}
1.6&-0.9&0.7&1.2&-1.4&0.5\\
-1.3&1.5&-0.8&0.6&1.0&-1.1\\
0.8&1.1&1.4&-1.5&-0.5&0.9\\
-0.7&-1.6&1.0&1.4&0.8&-0.6\\
1.2&0.6&-1.5&-0.9&1.3&1.0
\end{bmatrix},\quad c=(1,2,5,2,3),
\end{equation}
let $r_i=\lVert z-C_i\rVert_2^2$ and $C_\Sigma=13$.  The two normalized
objectives are
\begin{align}
f_1(z)&=\frac{C_\Sigma-
 \sum_{i=1}^{5}c_i\exp(-r_i/\pi)\cos(\pi r_i)}{2C_\Sigma},
 \label{eq:langermann}\\
f_2(z)&=\frac{-20\exp[-0.2\sqrt{\frac16\sum_j z_j^2}]
 -\exp[\frac16\sum_j\cos(2\pi z_j)]+20+e}{20+e-e^{-1}}.
 \label{eq:ackley}
\end{align}
$f_1$ is a bounded Langermann-style landscape and $f_2$ is normalized Ackley.
Their optima conflict: Ackley is minimized at $z=0$, whereas the Langermann
centers are displaced.  The reference point is $(1.1,1.1)$.  Nominal dimension
is high but effective dimension is six; the benchmark tests whether the
surrogate can recover this dense projection under a strongly nonlinear outer
map.

\subsection{Shallow-lake management: $d=100$, $K=21$}

Each $x_t$ controls annual anthropogenic phosphorus release
$u_t=0.1x_t$ over 100 years.  The lake is simulated under 100 fixed
common-random-number inflow scenarios.  For scenario $s$, $P_0^{(s)}=0$ and
for $t=1,\ldots,99$,
\begin{equation}
P_t^{(s)}=0.58P_{t-1}^{(s)}+
\frac{(P_{t-1}^{(s)})^2}{1+(P_{t-1}^{(s)})^2}
+u_{t-1}+\xi_{t-1}^{(s)},
\label{eq:lake-dynamics}
\end{equation}
where
$\xi_t^{(s)}\sim\operatorname{LogNormal}(-3.52,0.105^2)$ was drawn once with
seed 20,240,820.  Fixing scenarios makes the oracle deterministic while
retaining uncertainty in the expected physical response.  Let
$\bar P_t=100^{-1}\sum_sP_t^{(s)}$.  The first 20 components are five-year
block peaks,
\begin{equation}
h_b=\frac1{2.4}\max_{t=5(b-1),\ldots,5b-1}\bar P_t,
\quad b=1,\ldots,20.
\end{equation}
The final component is normalized discounted utility,
\begin{equation}
h_{21}=\frac{0.4\sum_{t=0}^{99}0.98^tu_t}
 {0.4(0.1)\sum_{t=0}^{99}0.98^t}.
\end{equation}
The minimized objectives and reference point are
\begin{equation}
 f_1=\max_{b=1,\ldots,20}h_b,\qquad
 f_2=1-h_{21},\qquad r=(1.05,1.05).
 \label{eq:lake-objectives}
\end{equation}
This is a deterministic two-objective slice of the shallow-lake dynamics of
Carpenter et al.\ \cite{carpenter1999lake} as used in multi-objective robustness
studies \cite{ward2015tipping}; it is not the original four-objective stochastic
benchmark.  The last decision $u_{99}$ contributes to utility but, following
the implemented indexing, not to the 100 recorded lake states.  The composite
representation retains which five-year period controls water quality before
the outer maximum discards that information.

\section{Optimization configurations and runtime}
\label{app:runtime}

The implementation uses double-precision PyTorch tensors with BoTorch and
GPyTorch models \cite{balandat2020botorch}.  Standard and spherical models are
fit with \texttt{fit\_gpytorch\_mll} using exact or summed marginal log
likelihoods; configured kernel and likelihood priors are included in that
fitting objective.  Table~\ref{tab:optimization-config} collects the proposal
settings that determine the evaluation accounting.  Additional kernel and
trust-region details are given in Appendix~\ref{app:implementation}.

\begin{table}[ht]
\centering
\caption{Optimization configurations used for the reported experiments.}
\label{tab:optimization-config}
\footnotesize
\setlength{\tabcolsep}{4pt}
\begin{tabular}{lcccl}
\toprule
Method family & Initial & Adaptive & Proposal structure & Raw/restarts \\
\midrule
qLogEHVI & 5 & 40 & sequential & 128/8 \\
Low-dimensional STCH & 5 & 40 & 5 weights $\times$ 8 & 128/8 \\
Spherical STCH & 20 & 100 & 5 weights $\times$ 20 & 256/10 \\
MORBO & 20 & 100 & 20 batches $\times$ 5 & 512/-- \\
\bottomrule
\end{tabular}
\end{table}

\begin{table}[ht]
\centering
\caption{Relative composite gain in time-averaged hypervolume over the adaptive
phase, and mean composite/direct solver-runtime ratio.  Runtime includes GP
fitting and acquisition but excludes one-time benchmark construction.}
\small
\setlength{\tabcolsep}{5pt}
\begin{tabular}{llrr}
\toprule
Benchmark & Method family & HV-curve gain & Runtime ratio \\
\midrule
Color formulation & qLogEHVI & $+14.87\%$ & $4.45\times$ \\
 & STCH & $+2.70\%$ & $2.48\times$ \\
DTLZ2 & qLogEHVI & $+119.25\%$ & $2.08\times$ \\
 & STCH & $+67.19\%$ & $1.37\times$ \\
Four-bar truss & qLogEHVI & $+2.91\%$ & $2.85\times$ \\
 & STCH & $+5.16\%$ & $3.07\times$ \\
Langermann--Ackley & spherical STCH & $+6.67\%$ & $6.47\times$ \\
 & MORBO & $+3.37\%$ & $3.84\times$ \\
Shallow lake & spherical STCH & $+4.80\%$ & $10.49\times$ \\
 & MORBO & $+2.59\%$ & $5.93\times$ \\
\bottomrule
\end{tabular}
\end{table}

Mean runtimes in seconds (direct $\rightarrow$ composite) were:
color qLogEHVI $114.7\rightarrow511.0$, color STCH
$92.9\rightarrow230.9$; DTLZ2 qLogEHVI $79.3\rightarrow165.2$, DTLZ2 STCH
$99.7\rightarrow136.4$; truss qLogEHVI $83.1\rightarrow237.1$, truss STCH
$67.1\rightarrow206.1$; Langermann--Ackley spherical STCH
$1860.4\rightarrow12036.4$, Langermann--Ackley MORBO
$1104.9\rightarrow4241.9$; shallow-lake spherical STCH
$1808.7\rightarrow18967.7$, and shallow-lake MORBO
$227.6\rightarrow1349.3$.  These are descriptive cluster wall times and were
not normalized for node contention or hardware variation.

\end{document}
